% 123456_789012_U04_A2b.tex
\documentclass{article}
\usepackage[utf8]{inputenc}
\usepackage{amsmath}
\usepackage{geometry}
\geometry{a4paper, margin=1in}

\title{CP\_II - Blatt 2}
\author{Julia Els, 454291 - Felix Krueckel, 454343}
\date{30. April 2025}

\begin{document}
\maketitle

\section*{Herleitung der t-Statistik für gleichgroße Stichproben}

Gegeben ist die allgemeine Formel für die t-Statistik für zwei unabhängige Stichproben mit Stichprobengrößen $N_1$ und $N_2$ und unbekannter, aber als gleich angenommener Standardabweichung $\sigma$:
\begin{equation}
t = \frac{\bar{x}_1 - \bar{x}_2}{S_p \sqrt{1/N_1 + 1/N_2}}
\end{equation}
wobei $S_p$ die gepoolte Standardabweichung ist, berechnet aus der gepoolten Varianz $S_p^2$:
\begin{equation}
S_p^2 = \frac{(N_1 - 1) s_1^2 + (N_2 - 1) s_2^2}{N_1 + N_2 - 2}
\end{equation}
Hier sind $\bar{x}_1, \bar{x}_2$ die Stichprobenmittelwerte und $s_1^2, s_2^2$ die Stichprobenvarianzen.

Wir sollen die Formel für den Fall herleiten, dass die Stichprobengrößen gleich sind, d.h. $N_1 = N_2 = N$.

Setzen wir $N_1 = N$ und $N_2 = N$ in die Formel für die gepoolte Varianz $S_p^2$ ein:
\begin{align*}
S_p^2 &= \frac{(N - 1) s_1^2 + (N - 1) s_2^2}{N + N - 2} \\
&= \frac{(N - 1) (s_1^2 + s_2^2)}{2N - 2} \\
&= \frac{(N - 1) (s_1^2 + s_2^2)}{2(N - 1)}
\intertext{Da $N > 1$ für eine Varianzberechnung, können wir $(N-1)$ kürzen:}
S_p^2 &= \frac{s_1^2 + s_2^2}{2}
\end{align*}

Nun setzen wir $N_1 = N$, $N_2 = N$ und das Ergebnis für $S_p^2$ in die Formel für $t$ ein. Der Nenner des t-Wertes ist $S_p \sqrt{1/N_1 + 1/N_2}$:
\begin{align*}
\text{Nenner} &= \sqrt{S_p^2} \sqrt{\frac{1}{N} + \frac{1}{N}} \\
&= \sqrt{\frac{s_1^2 + s_2^2}{2}} \sqrt{\frac{2}{N}} \\
&= \sqrt{\frac{(s_1^2 + s_2^2)}{2} \cdot \frac{2}{N}} \\
&= \sqrt{\frac{s_1^2 + s_2^2}{N}}
\end{align*}

Somit ergibt sich für die t-Statistik im Fall $N_1 = N_2 = N$:
\begin{equation}
t = \frac{\bar{x}_1 - \bar{x}_2}{\sqrt{\frac{s_1^2 + s_2^2}{N}}}
\end{equation}
Dies ist die aus der Vorlesung bekannte Formel für den t-Test bei zwei gleichgroßen, unabhängigen Stichproben mit unbekannter, aber gleicher Varianz.

\end{document}